\section{Formal verification and reproducible supplements}
\label{sec:formal-overview}

The accompanying Lean development verifies specified parts of the
paper. It has 64 modules with 909 owned declarations, plus five imported
local support modules. The portable sources, pinned Lean~4.33.1
toolchain and Mathlib dependency manifest, declaration coverage maps,
and a targeted verification runner are in \path{supplement/lean/}. A
separate formal supplement relates the theorem interfaces to the
mathematical statements. The table lists the five packages; sizes are
modules / owned declarations, and package numbers only label the
supplied audit partitions.

\begin{center}
\small
\begin{tabular}{@{}p{0.09\textwidth}p{0.12\textwidth}p{0.69\textwidth}@{}}
\hline
Package & Size & Verified mathematical scope\\
\hline
27 & 11 / 178 & Theorem~\ref{thm:certificate}, its HHC specialization,
AHC sequence equivalence, and the principal separation and cone lemmas.\\[3pt]
28 & 10 / 158 & Closed-hull properness, the unconditional whole-space Shor
criterion and its AHC consequences, exact signed-coordinate SDP tests,
Dines' theorem, Example~\ref{ex:strip}, and the strict-feasibility
counterexample of Example~\ref{ex:closed} ($S=\varnothing$, the form of
$T$, proper hull, trivial convex certificates, and HHC); its interior
and good-multiplier arguments are not formalized.\\[3pt]
29 & 18 / 248 & Actual HHC for $S_r$, $r\geq2$, the eigenvalue-based good
cone, indispensable uncountably many strict rays, and impossibility of a
finite weak good-aggregation description.\\[3pt]
30 & 10 / 93 & Exact hull and lift assertions of Theorem~\ref{thm:ball-hulls}: the two-point
construction in every $r\geq2$, exact strict and closed hulls, PD/PSD lifts,
weak-system hull equality, and all-good intersection formulas.\\[3pt]
31 & 15 / 232 & Explicit Euclidean approximation bounds, the
$\Theta(N^{-2})$ rate, tolerance families, and exact rational meshes with
logarithmic coefficient sizes.\\
\hline
\end{tabular}
\end{center}

Two formal statements differ from the paper. Package~28 checks
$n(n+1)+2n$ signed coordinate objectives; the shorter $2n+1$
trace-and-vector test of Corollary~\ref{cor:finite-sdp} is proved only
in the paper. Package~31 proves the upper bound of
Theorem~\ref{thm:approximation-rate} and the lower bound
$1/(2000N^2)$, which implies the weaker bound
$\sqrt2(\log2)^2/(1600N^2)$ stated in that package; the lower constant $\sqrt2/2000$ used
in this paper is not formalized. The formal approximation model
explicitly transports the Euclidean norm on all $2r$ coordinates rather
than using a product maximum norm.

The formal definitions use actual feasible sets, ordinary convex hulls,
nonnegative aggregation weights, and homogeneous eigenvalues. HHC of the
example is proved, not assumed; the exact-hull package proves its
two-point construction without a general HHC hull theorem as a premise;
and the certificate proof derives its separation weights and nonzero
coefficient limit. None of these enters the theorem interfaces as a
hidden assumption.

For each package, warning-free targeted compilation, a transitive axiom
audit of all owned declarations, and per-module kernel replay have passed
with recorded source fingerprints. The audit permits only
\texttt{propext}, \texttt{Classical.choice}, and \texttt{Quot.sound}.
The supplied runner reproduces these checks for the portable project and
records whether its inputs are unchanged. Kernel replay uses Lean's
kernel and the imported dependencies; it does not rebuild and replay all
dependencies, and it is not an independent proof-assistant
implementation. Exact rational checks of the examples and the
figure-generation program are also supplied; they check stated
identities and witnesses, not universal theorems.

The Lean verification does not cover, among others, the general sharp
Gram-map theorem,
the arbitrary-quadratic obstruction, countable dense weak sufficiency,
the four-aggregation PDLC theorem,
the many-row appendix, the local-certificate application, the
single-objective duality result, the sufficient criteria of
Proposition~\ref{prop:stable}, the general closure equality of
Proposition~\ref{prop:shor-closure}, or
Examples~\ref{ex:nonclosed-shor}, \ref{ex:ahc-not-hhc},
\ref{ex:hhc-not-stable}, and~\ref{ex:ordinary-hc}. The spectral quantitative alternative
to the certificate proof is not fully formalized. Literature priority
and numerical solver correctness are outside every package. For these
results the paper supplies the mathematical proofs and source
comparisons.
